\documentclass[11pt]{article}
\usepackage{fontspec}
\usepackage{amsmath,amssymb}
\setmainfont{DejaVu Serif}
\usepackage[margin=1in]{geometry}
\title{Switches instead of gates:\\
does a contact model augment the TSR charges?}
\author{Grok, identity j15marti\\ repository \texttt{j15marti/pnp}, branch \texttt{main}}
\date{3 October 2026, 10:08 PT}
\begin{document}
\maketitle
% GROK SWITCH1 :: 2026-10-03_1008_GROK_SWITCH1

P versus NP remains open.

This note is a literature pass on one proposal, made 2026-10-03 at about 10:08 PT: model the problem as individual on/off switches, not as AND/OR/NOT gates, and ask whether that augments the time--state--representation architecture already committed.
It does not edit root \texttt{main.tex}.
It does not add a workflow.
It does not rerun \(D(4,8)\).
It does not claim that P differs from NP, and it does not claim that P equals NP.

The architecture is left as written in the three files below.
This note only asks what the opened literature does to its charges.
Those charges are: dynamic time \(T\) (longest source-to-answer path), mutable working memory \(S\), stored representation \(R\), continuation nonmergeability (two histories that some common continuation separates cannot share one deterministic state), and the open compiler from an arbitrary SAT decider into the charged model.

\section{Architecture pointed at, not rewritten}
Opened on this pass, branch \texttt{main}:
\begin{itemize}
\item \texttt{research/agents/grok/2026-10-03\_0948\_GROK\_TSR1.tex}, blob \texttt{87f66829798744adb35bae1cc1c003bfcf0a91c1}.
It proposes the Pareto set of \((T,S,R)\), defines continuation incompatibility, and leaves pebbling and the compiler \textbf{OPEN}.
\item \texttt{research/agents/grok/2026-10-03\_0950\_GROK\_TSR\_ITER1.tex}, blob \texttt{c498c2a267c3e7bac6aa25bf2faff8a7f41c2557} as read on this pass.
The witnessed ordered DAG for \(D(4,8)\) has \(720\) nonterminal nodes, longest decision path \(T=31\), level-\(16\) width \(66\), and the cell \((t,k)=(16,4)\) tight at \(K=32\).
\item \texttt{research/agents/grok/2026-10-03\_1000\_GROK\_TSR\_ITER2.tex}, blob \texttt{741af07e4ba205f6c10028f81b58f548cc802956}.
Forgetting the whole mask at that cut drops the stored configurations there from \(66\) to \(5\), and the exhibited repair raises \(T\) from \(31\) to \(47\) and \(R\) from \(720\) to \(1330\).
That point is not a \((T,R)\) Pareto improvement.
ITER2 does not invoke a pebble-game theorem.
\end{itemize}
Nothing above is recomputed here.
The switch question is whether those charges get a new theorem by reading the same DAG as contacts.

\section{What was opened}
Each item below was opened in this pass (PDF text or a rendered page).
Lupanov's 1958 papers and Reingold's 2008 paper were not opened.
Where a theorem of theirs is used, the wording is the wording of a book or survey that was opened, and the status line says so.
No paper, year, theorem, or number was filled in from memory.

\subsection*{Shannon, two-terminal switching circuits}
Opened: Claude E.\ Shannon, ``The Synthesis of Two-Terminal Switching Circuits,'' \emph{Bell System Technical Journal} 28 (1949), 59--98, the scanned PDF.
The introduction, as printed, interprets a symbol \(X\) as a make contact on a relay or switch and \(X'\) as a break contact.
Addition is series connection of those contacts.
Multiplication is parallel connection.
The constants \(0\) and \(1\) are a permanently closed circuit and a permanently open circuit.
That is the on/off switch model, not a gate model.
Theorem~7, as read from the journal pages in that scan, states that for all sufficiently large \(n\), all functions of \(n\) variables except a fraction \(\delta\) require at least \((1-\varepsilon)2^n/n\) elements, where \(\varepsilon\) and \(\delta\) are arbitrarily small positive numbers, and hence \(\lambda(n)>(1-\varepsilon)2^n/n\).
The argument is counting: there are not enough networks with fewer branches, once literals are assigned to branches, to realize every function.
Shannon writes that \(\lambda(n)\) is of order \(2^n/n\).
He also writes that he has not proved the stronger opinion that almost all functions have complexity greater than \(1-\varepsilon\); the network count in his lemma was not tight enough.
The scan's upper-bound constants in Theorem~6 were not clean enough to quote, so they are not quoted.
Status of Theorem~7: \textbf{KNOWN THEOREM} about two-terminal contact networks.
It is a counting fact about almost all Boolean functions.
It is not a lower bound for SAT, and it is not a lower bound for \(D(4,8)\).

\subsection*{Jukna, \emph{Boolean Function Complexity}}
Opened: Stasys Jukna, \emph{Boolean Function Complexity: Advances and Frontiers}, the draft PDF \texttt{bool-V7.pdf} (file title page marked ``Draft'').
Section~1.3 defines a switching network, also called a contact scheme, as an undirected graph with source and target, wires labeled by literals.
An input switches each labeled wire on or off.
The network accepts when some source-to-target path is entirely on.
Unlabeled wires may be contracted.
A parallel-serial network (\(\pi\)-scheme) is one whose wires can be oriented so that every source-to-target path becomes a directed path.
Proposition~1.10 states that every DeMorgan formula can be simulated by a \(\pi\)-scheme of the same size, and vice versa.
Series-parallel switches are formulas.
They are not a new model past formulas.

The same section states \(\mathrm{NBP}(f)\le S(f)\le\mathrm{BP}(f)\), where \(S\) is undirected contact size, \(\mathrm{NBP}\) is nondeterministic branching-program size (directed contacts, unlabeled rectifiers allowed), and \(\mathrm{BP}\) is deterministic branching-program size.
It then states that an important result of Reingold (2008) translates to \(\mathrm{BP}(f_n)\le S(f_n)^{O(1)}\), and it calls that a ``P\,=\,NP type result for branching programs.''
Reingold's paper was not opened.
The polynomial comparison, as a stated theorem in the book that was opened, is \textbf{KNOWN THEOREM} in that text.
This note does not re-prove it.

The same book states Lupanov's bounds, again without this pass opening Lupanov.
Theorem~1.15 (Lupanov 1958a), as printed: every Boolean function \(f\) of \(n\) variables satisfies \(C(f)\le(1+\alpha_n)2^n/n\) with \(\alpha_n=O((\log n)/n)\), for circuit size \(C\).
Section~1.4.4 (Lupanov 1958b), as printed: the switching-network Shannon function satisfies
\[
\Bigl(1+\frac{2\log n-O(1)}{n}\Bigr)\frac{2^n}{n}
\le S(n)
\le\Bigl(1+\frac{O(1)}{\sqrt{n}}\Bigr)\frac{2^n}{n}.
\]
The same section attributes to Shannon (1949) the coarser two-sided bound \((1-\varepsilon)2^n/n<S(n)<2^{n+3}/n\).
That upper constant is Jukna's attribution.
It was not re-derived from the Shannon scan.
For formulas, the book states that some functions have formulas at least \(n/\log n\) times larger than circuits and switching networks, whose Shannon functions are of order \(2^n/n\).

Markov (1962), as printed in the same section: for a monotone function, monotone nondeterministic branching-program size is at least the product of minimum minterm length and minimum maxterm length, and for the threshold function this product is tight, \(\mathrm{NBP}_+(\mathrm{Th}_n^k)=k(n-k+1)\).
In particular the book prints \(\mathrm{NBP}_+(\mathrm{Th}_n^2)=2n-3\).
It also prints that monotone undirected contact size satisfies \(S_+(\mathrm{Th}_n^2)=\Omega(n\log n)\), with a matching \(n\log_2 n\) upper bound for \(n\) a power of two even in the class of \(\pi\)-schemes, while non-monotone undirected contacts satisfy \(S(\mathrm{Th}_n^2)\le 3n-4\).
Direction is not free: for this monotone function the directed model is linear and the undirected monotone model is larger by a logarithmic factor.
Status of these printed statements: \textbf{KNOWN THEOREM} as theorems of the opened draft.
They are not theorems about \(D(4,8)\).

\subsection*{Razborov, FCT 1991}
Opened: Alexander A.\ Razborov, ``Lower Bounds for Deterministic and Nondeterministic Branching Programs,'' the PDF at \texttt{people.cs.uchicago.edu/\textasciitilde razborov/files/fct.pdf}, including the rendered introduction page that carries the Ne\v{c}iporuk display.
A switching-and-rectifier network is directed.
Edges are labeled \(x_i=0\) or \(x_i=1\), or left free.
Size \(\mathrm{RS}(f)\) counts labeled edges.
A switching network is the undirected version; free edges do not help, because they contract.
Size \(S(f)\) counts edges.
A branching program is the deterministic special case: out-degree \(2\), the two out-edges testing one variable both ways.
The extracted text prints a chain under a relation, written \(\preceq\) here, which that text defines by \(F\preceq G\) meaning \(F=O(G)\):
\[
C(f)^{1/3}\preceq\mathrm{RS}(f)\preceq S(f)\le\mathrm{BP}(f)\preceq L(f),
\]
where \(C\) is circuit size and \(L\) is formula size over \(\{\neg,\wedge,\vee\}\).
He states \(\mathrm{LOGSPACE}/\mathrm{poly}=\mathrm{BP}(\mathrm{poly})\) and \(\mathrm{NLOGSPACE}/\mathrm{poly}=\mathrm{RS}(\mathrm{poly})\).
For undirected switching networks he argues \(\mathrm{LOGSPACE}/\mathrm{poly}=S(\mathrm{poly})\) because undirected accessibility is in random logarithmic space, and he concludes \(\mathrm{BP}(f)\le S(f)^{O(1)}\).
He writes that no constructive proof of that last simulation was known.
He also prints \(\mathrm{RS}(f)\preceq\mathrm{NBP}(f)\le\mathrm{RS}(f)^{O(1)}\) for his guessing-node variant of nondeterministic branching programs.
The introduction, as rendered, states that the best bound then known for an NP function, and actually holding for a function in P, is Ne\v{c}iporuk's 1966 bound of order \(\Omega\bigl(n^2/(\log n)^2\bigr)\) for switching networks, the same bound for branching programs, and order \(\Omega\bigl(n^{3/2}/\log n\bigr)\) for switching-and-rectifier networks, and that superpolynomial lower bounds for NP functions were still open in all three models.
He also writes that circuits capture time, while these network models capture space.
Status: \textbf{KNOWN THEOREM} as stated in the opened survey.
Not a separation of P from NP.
Not a bound on \(D(4,8)\).

\subsection*{Wegener, \emph{The Complexity of Boolean Functions}}
Opened: Ingo Wegener, \emph{The Complexity of Boolean Functions}, Wiley--Teubner, 1987, the PDF facsimile, Chapter~14, book pages 414--417.
Definition~1.1 is a deterministic branching program: one variable test per inner node, two successors.
Definition~1.2: size is the number of inner nodes, depth is the longest path, and depth is called a measure of computation time.
Theorem~1.1 (Cobham 1966; Pudl\'ak and Z\'ak 1983), as printed: nonuniform space \(S(f_n)\) is \(O(\log\max\{\mathrm{BP}(f_n),n\})\), and \(\mathrm{BP}(f_n)=2^{O(h(n))}\) where \(h(n)=\max\{S(f_n),\log n\}\).
Theorem~1.2, as printed: \(C(f)\le 3\,\mathrm{BP}(f)\), depth at most twice branching-program depth, and \(\mathrm{BP}(f)\le L(f)+1\) over \(\{\wedge,\vee,\neg\}\).
Theorem~1.3, as printed: \(\mathrm{BP}(f)\ge(1/3)\,2^n/n\) for almost all \(f\in B_n\), and \(\mathrm{BP}(f)=O(2^n/n)\) for every \(f\in B_n\).
The chapter that was opened compares branching programs with circuits, formulas, and nonuniform Turing space.
It does not contain a proof that undirected contact size and branching-program size are polynomially related.
Contact-circuit papers appear in the bibliography of the book (Hansel, Krichevskii, Red'kin).
They were not opened.
The expectation that this 1987 book would itself contain the contact-network comparison was not met by the chapter that was read.
Status of Theorems~1.1--1.3: \textbf{KNOWN THEOREM}.
They do not mention \(D(n,m)\) and they do not separate P from NP.

\subsection*{Jukna, null-chains}
Opened: Stasys P.\ Jukna, ``The Effect of Null-Chains on the Complexity of Contact Schemes,'' the PDF of the FCT 1989 paper.
In that paper a contact scheme is a labeled directed graph.
Size is the number of edges.
A chain is a source-to-output path.
A null-chain is a chain that contains some variable and its negation, so its conductivity is identically zero.
The abstract states that the main tool for reducing scheme size is the use of null-chains, because those chains let one merge non-isomorphic subschemes, and that no simple function, say in NP, is known to require contact-scheme size \(\sim(n^3)\), in the symbol printed there.
The abstract states that schemes without null-chains (Theorems~1--2), schemes with restricted null-chain topology (Theorem~3), and schemes with a restricted number or length of null-chains (Theorem~4) all receive nearly-exponential lower bounds, and that null-chains do not help when the scheme must realize sufficiently many prime implicants (Theorem~5).
The explicit inequalities in the body did not survive text extraction cleanly, so they are not copied.
One-time-only branching programs are identified in the paper as a special case of contact schemes without null-chains.
Status of the abstract claims: \textbf{KNOWN THEOREM} as claims of the opened paper.
The \(\sim(n^3)\) symbol is his, not a conversion.
The paper says the unrestricted contact-scheme lower-bound problem for NP functions was unsolved.

\section{What this does to the five charges}
The literature that was opened says the following, in plain words.
Undirected contact size and deterministic branching-program size are polynomially related: \(S(f)\le\mathrm{BP}(f)\) and \(\mathrm{BP}(f_n)\le S(f_n)^{O(1)}\), the second as Jukna states it via Reingold, and as Razborov stated it in 1991 with no constructive proof then known.
Directed contacts sit with nondeterministic branching programs and with nonuniform nondeterministic logspace, again up to a polynomial.
Series-parallel contacts are formulas, exactly, at the same size.

That polynomial relation is the whole answer to the proposal.
Reading variables as on/off switches changes constants and changes the natural sentence for nonmergeability.
It does not by itself evade the compiler gap, and it does not close it.

\subsection*{\(T\), path length}
Wegener's opened chapter charges branching-program depth as time: one variable test, then one successor.
A contact scheme accepts by existence of an on-path, not by a single deterministic walk.
The opened comparisons do not produce a shorter path measure that escapes branching-program depth.
Jukna's polynomial simulation says an undirected contact scheme of size \(S\) has some deterministic branching program of size \(S^{O(1)}\).
That program's depth is a \(T\) for the same function, and nothing opened bounds that depth by a smaller function of our \(T=31\).
Relabeling the committed \(D(4,8)\) DAG as contacts does not change the exhibited longest decision path.
This note does not recompute it.
Switches do not augment \(T\).

\subsection*{\(S\), live mutable state}
In the deterministic branching program the live state is the current node.
Wegener's Theorem~1.1 says the number of such nodes and nonuniform space determine each other in the usual exponential way.
For undirected contacts, the polynomial simulation puts the same state bound, up to a polynomial in the representation size, on the deterministic side.
Directed rectifiers are nondeterministic space, which is a different class, not a smaller working memory.
The level-\(16\) width \(66\), and the \(32\) mutually separating histories inside \(k=4\), are deterministic state counts already.
A switch reading does not split them, and it does not force a new bit of mutable memory beyond the branching-program reading ITER1 already used.
Switches do not augment \(S\).

\subsection*{\(R\), representation size}
For almost all Boolean functions, Jukna's statement of Lupanov says both circuit size and undirected contact size are \((1+o(1))2^n/n\).
The switch syntax does not reduce worst-case representation below gates.
Series-parallel switches are formulas, and the formula Shannon function in the same book is on the order of \(2^n/\log n\), which is worse.
The explicit Ne\v{c}iporuk bound Razborov quotes for switching networks is \(\Omega\bigl(n^2/(\log n)^2\bigr)\), the same order he quotes for branching programs, and not superpolynomial.
Jukna's 1989 abstract still had no NP function at contact size \(\sim(n^3)\).
None of this is a new representation lower bound for the \(720\)-node DAG, and none of it is a lower bound for every SAT algorithm.
Switches do not augment \(R\).

\subsection*{Continuation nonmergeability}
This is the only place a contact-specific sentence is easier than the gate sentence.
Call it a \textbf{CANDIDATE} augmentation, not a proof.
Jukna 1989: a null-chain conducts for no input, and its role is to merge subschemes that are not isomorphic.
That is the contact form of sharing a vertex across two histories that a legal continuation would have separated.
The deterministic charge in TSR1 forbids that share: one state plus one continuation cannot realize two values.
The opened lower bounds that are nearly exponential are proved when null-chains are forbidden or restricted, that is, when the merge is not available.
When the merge is available, the same abstract says the NP lower-bound problem is still open below \(\sim(n^3)\).
So the switch model names the glue.
It does not show that forbidding the glue on our multisign histories costs more than the branching-program width already counted at cell \((16,4)\).
It does not show that every correct machine must keep those histories unmerged in its stored representation.
Status of this reading: \textbf{CANDIDATE}.
It is not \textbf{PROVED/EXACT} for the multisign DAG, and it is not a theorem about SAT.

\subsection*{The compiler}
The load-bearing arrow in TSR1 stays \textbf{OPEN}.
Razborov writes that circuit computations capture time in Turing computation, and that switching networks and branching programs capture space.
The class inclusions he prints from that comparison are \(\mathrm{P}\supseteq\mathrm{NLOGSPACE}\supseteq\mathrm{LOGSPACE}\).
This pass did not open a proof that every polynomial-time machine has polynomial circuits, so that simulation is not claimed as a theorem of this note.
What the opened texts do say is enough for the negative point.
Polynomial undirected contact schemes are, by combining Razborov's \(\mathrm{BP}(\mathrm{poly})=\mathrm{LOGSPACE}/\mathrm{poly}\) with the polynomial comparison \(S\) versus \(\mathrm{BP}\) stated above, nonuniform logspace.
Polynomial directed contact schemes are, by his equality \(\mathrm{RS}(\mathrm{poly})=\mathrm{NLOGSPACE}/\mathrm{poly}\), nonuniform nondeterministic logspace.
Neither statement says that every polynomial-time SAT algorithm is a small switching network.
Moving from gates to switches moves into the space models he describes, which sit inside the inclusions above, not outside them.
The switch syntax therefore does not supply the missing compiler.
If a future proof showed that every SAT algorithm gives a small contact scheme, that would be a new theorem about space, past anything opened here, and this note does not claim it.

Pebbling is not this augmentation.
TSR1 leaves the pebbling inequality \textbf{OPEN}.
ITER2 charges one concrete forget-and-reread and does not call it a pebble theorem.
A contact, in the papers opened, is an edge that the input turns on or off.
A pebble, in the architecture's own wording, is a token the computation places and later removes.
Those are different objects.
Calling the existing pebbling charge ``switch-like'' does not turn the open pebbling arrow into a theorem, and it does not change the ITER2 point \((T,R)=(47,1330)\).

\section{Plain answer}
The on/off switch model does not augment the approach.
Undirected contacts are polynomially the same as deterministic branching programs.
Directed contacts are polynomially the same as nondeterministic branching programs.
Series-parallel contacts are formulas.
Worst-case representation stays on the Shannon--Lupanov order \(2^n/n\), the same order as circuits, and the explicit NP bounds that were opened stay far from superpolynomial.
What the switch language does do is name continuation sharing as a null-chain: a dead path used only to glue non-isomorphic subschemes.
That name is a candidate description of nonmergeability.
The strong lower bounds in the opened null-chain paper hold when that glue is forbidden, which is the restriction already present in a deterministic state, and they are not known for unrestricted contact schemes.
Nothing opened compiles an arbitrary SAT decider into the charged \((T,S,R)\) model.
That compiler remains open.
P versus NP remains open.

\section{Status ledger}
\begin{itemize}
\item P versus NP: \textbf{OPEN}.
\item Shannon's Theorem~7, as read from the 1949 paper: \textbf{KNOWN THEOREM}. Counting lower bound for contact networks. Not SAT.
\item Jukna's printed comparisons \(\mathrm{NBP}\le S\le\mathrm{BP}\), the Reingold translation \(\mathrm{BP}\le S^{O(1)}\), Proposition~1.10, Markov's threshold bound, and the Lupanov Shannon functions as printed in the opened draft: \textbf{KNOWN THEOREM} in that draft. Lupanov's originals and Reingold's paper were not opened.
\item Razborov's chain of inequalities, the logspace class equalities, and the Ne\v{c}iporuk orders as printed in the opened survey: \textbf{KNOWN THEOREM}. He records that the undirected-to-branching-program simulation had no constructive proof in that paper.
\item Wegener, Theorems~1.1--1.3, as printed on the opened pages: \textbf{KNOWN THEOREM} for branching programs, circuits, formulas, and nonuniform space.
\item Jukna 1989, null-chains as the merge tool, and nearly-exponential bounds when they are restricted: \textbf{KNOWN THEOREM} as stated in the opened abstract. The explicit body inequalities were not copied.
\item Reading the committed \(D(4,8)\) DAG as contacts: \textbf{CANDIDATE}. Not a new computation.
\item Null-chains as the contact form of continuation nonmergeability: \textbf{CANDIDATE} augmentation. Not a proof.
\item Compiler from an arbitrary SAT decider into the charged model: \textbf{OPEN}.
\item Pebbling price: \textbf{OPEN}, unchanged.
\item Nothing in this pass is \textbf{FALSIFIED/RETIRED}.
\end{itemize}

% END_OF_VERIFIED_REPORT::GROK::2026-10-03_1008_GROK_SWITCH1
\end{document}
